\documentclass{antiquebook}
\SetOrnamentFont{CenturymodernTT-Regular.otf}

\usepackage{lipsum}
\usepackage[math]{blindtext}

\author{L.D. Landau \& E.M. Lifshitz}
\title{{Kurz teoretické fyziky}}

%\input{glossary.tex}

\begin{document}
	\frontmatter
	\maketitle
	\tableofcontents
	\mainmatter
	\pagestyle{fancy}

\chapter{Základní koncepty kvantové mechaniky}
\section{Princip neurčitosti}
Pokusíme-li se vysvětlit atomové jevy prostřednictvím klasické mechaniky a elektrodynamiky, získáváme výsledky v očividném nesouladu s experimentem. Jeden z nich obdržíme aplikací obyčejné elektrodynamiky na model atomu, v němž se elektrony pohybují okolo jádra na klasických obězných drahách. Při takovémto pohybu by měly elektrony spojitě vyzařovat elektromagnetické vlny, jakožto při jakémkoliv zrychleném pohybu. Tímto by však elektrony postupně ztratily svou energii a spadly do jádra. Dle klasické elektrodynamiky je tedy atom nestabilní, což je spor s realitou.\par
Tento rozpor mezi teorií a experimentem ukazuje, že konstrukce teorie aplikovatelné na atomové jevy--to jest jevy částic velmi malých hmotností na velmi malých vzdálenostech--nutně vyžaduje fundamentální modifikace základních fyzikálních konceptů a zákonů.\par
Výchozím bodem pro zkoumání těchto modifikací je vhodné uvažovat experimentálně pozorovaný fenomén zvaný \emph{elektronová difrakce}.{\renewcommand{\thefootnote}{\fnsymbol{footnote}}\footnote[2]{Fenomén elektronové difrakce byl ve skutečnosti objeven až po vynálezu kvantové mechaniky. Pro naši diskusi však není klíčová historická posloupnost vývoje teorie, nýbrž snaha ji konstruovat takovým způsobem, aby byly patrné souvislosti mezi základními principy kvantové mechaniky a experimentálními pozorováními.}} Je pozorováno, že při průchodu homogenního svazku elektronů krystalem dochází ve výstupním svazku ke vzniku střídajících se maxim a minim intenzity, velmi podobně jako při difrakci elektromagnetických vln. Za určitých podmínek tedy chování hmotných částic--zde elektronů--vykazuje vlastnosti typické pro vlnové procesy.
\end{document}